%%%%%%%%%%%%%%%%%%%%%%%%%%%%%%%%%
%
%       EXERCÍCIO
%
%%%%%%%%%%%%%%%%%%%%%%%%%%%%%%%%%

\definevalorquestao

\begin{exercicioBanco}[\valorquestao]
Considere a função
%
\[
    \begin{array}{rcl}
        f:\bbR & \to & \bbR \\
        x & \mapsto &
        \left\{
        \begin{array}{cl}
            x+2, & \text{se } x \le -1, \\[2pt]
            1, & \text{se } -1 < x < 1, \\[2pt]
            -2x+3, & \text{se } x \ge 1.
        \end{array}
        \right.
    \end{array}
\]
%
Determine a soma dos zeros dessa função.

% Define as alternativas
\newcommand{\alternativas}{%
    \begin{center}
        \begin{tabularx}{\textwidth}{XXXXX}
            (a) \(-\dfrac{1}{2}\). &
            (b) \(-2\). &
            (c) \(0\). &
            (d) \(\dfrac{3}{2}\). &
            (e) \(1\).
        \end{tabularx}
    \end{center}
}

% Define a resposta correta
\newcommand{\resposta}{A}

% Lógica condicional para exibição
\ifdefstring{\atividade}{lista}{%
    \alternativas
    \vspace{0.5em}
    
    \noindent\textbf{Resposta:} letra \textbf{\resposta}.
}{%
    \ifdefstring{\modo}{objetivo}{%
        \alternativas
    }{%
        % modo = discursiva → não mostra alternativas
    }
}
\end{exercicioBanco}

